\chapter{IR and Compiler Pass Pipeline}

The compiler driver in \pathcode{src/einlang/compiler/driver.py} orchestrates
the implementation. It parses source, constructs the module system, merges
differential rules, resolves names, lowers AST to IR, runs the pass pipeline,
tree-shakes unused code, and returns a compilation result. The runtime receives
IR that has already been annotated with the facts needed by the backend.

\begin{figure}[p]
\centering
\resizebox{\textwidth}{!}{\begin{tikzpicture}[
  font=\sffamily\footnotesize,
  >=Latex,
  stage/.style={draw=figLine,line width=0.6pt,rounded corners=2pt,fill=white,align=left,inner sep=6pt,text=figInk,minimum width=2.9cm,minimum height=1.2cm},
  phase/.style={stage,fill=figBlueBg,draw=figBlue},
  analysis/.style={stage,fill=figGreenBg,draw=figGreen},
  lowering/.style={stage,fill=figGoldBg,draw=figGold},
  runtime/.style={stage,fill=figRoseBg,draw=figRose},
  side/.style={draw=figLine,line width=0.6pt,rounded corners=2pt,fill=figSlateBg,align=left,inner sep=6pt,text=figInk,minimum width=3.0cm},
  arrow/.style={->,line width=0.8pt,draw=figLine},
  bus/.style={line width=0.8pt,draw=figLine},
  rowlabel/.style={font=\sffamily\bfseries,text=figInk,align=right}
]

\node[rowlabel] at (-1.5,0) {Front end};
\node[phase] (parse) at (1.25,0) {
  \textbf{Parse and AST}\\
  Lark parser\\
  transformers\\
  source spans
};
\node[phase] (resolve) at (4.95,0) {
  \textbf{Identity}\\
  modules/imports\\
  DefIds\\
  scoped names
};
\node[phase] (ir) at (8.65,0) {
  \textbf{IR construction}\\
  slot nodes\\
  source locations\\
  DefId uses
};

\node[rowlabel] at (-1.5,-2.9) {Static\\analyses};
\node[analysis] (range) at (1.25,-2.9) {
  \textbf{Range}\\
  explicit domains\\
  implicit reads\\
  intersections
};
\node[analysis] (shape) at (4.95,-2.9) {
  \textbf{Shape}\\
  output axes\\
  array literals\\
  dependent ranges
};
\node[analysis] (type) at (8.65,-2.9) {
  \textbf{Type}\\
  primitives\\
  tensors\\
  calls and matches
};

\node[rowlabel] at (-1.5,-5.8) {High-level\\transforms};
\node[lowering] (mono) at (1.25,-5.8) {
  \textbf{Specialize}\\
  clone generics\\
  dynamic shapes\\
  incremental passes
};
\node[lowering] (ad) at (4.95,-5.8) {
  \textbf{Autodiff}\\
  rewrite \texttt{@}\\
  synthesize JVP/VJP\\
  lazy Jacobians
};
\node[lowering] (ein) at (8.65,-5.8) {
  \textbf{Tensor lowering}\\
  loops/guards\\
  reductions\\
  execution facts
};

\node[rowlabel] at (-1.5,-8.7) {Execution\\boundary};
\node[runtime] (rec) at (1.25,-8.7) {
  \textbf{Recurrence}\\
  same-step order\\
  window metadata\\
  circular buffers
};
\node[runtime] (validate) at (4.95,-8.7) {
  \textbf{Validate}\\
  casts\\
  exhaustiveness\\
  AD leak checks
};
\node[runtime] (backend) at (8.65,-8.7) {
  \textbf{Backends}\\
  NumPy visitor\\
  IREE subset\\
  profiling/debug
};

\node[side] (tcx) at (12.65,-1.45) {
  \textbf{TyCtxt}\\
  resolver state\\
  source files\\
  diagnostics\\
  analysis tables\\
  module loader
};
\node[side] (facts) at (12.65,-5.8) {
  \textbf{Attached facts}\\
  \texttt{range\_ir}\\
  \texttt{shape\_info}\\
  \texttt{type\_info}\\
  recurrence dims\\
  vectorization strategy
};

\draw[arrow] (parse.east) -- (resolve.west);
\draw[arrow] (resolve.east) -- (ir.west);
\coordinate (frontbus) at (8.65,-1.6);
\draw[bus] (ir.south) -- (frontbus);
\draw[bus] (frontbus) -- (1.25,-1.6);
\draw[arrow] (1.25,-1.6) -- (range.north);

\draw[arrow] (range.east) -- (shape.west);
\draw[arrow] (shape.east) -- (type.west);
\coordinate (analysisbus) at (8.65,-4.5);
\draw[bus] (type.south) -- (analysisbus);
\draw[bus] (analysisbus) -- (1.25,-4.5);
\draw[arrow] (1.25,-4.5) -- (mono.north);

\draw[arrow] (mono.east) -- (ad.west);
\draw[arrow] (ad.east) -- (ein.west);
\coordinate (loweringbus) at (8.65,-7.4);
\draw[bus] (ein.south) -- (loweringbus);
\draw[bus] (loweringbus) -- (1.25,-7.4);
\draw[arrow] (1.25,-7.4) -- (rec.north);

\draw[arrow] (rec.east) -- (validate.west);
\draw[arrow] (validate.east) -- (backend.west);

\draw[arrow] (tcx.west) -- ++(-0.55,0) |- (type.east);
\draw[arrow] (facts.west) -- (ein.east);
\draw[arrow] (facts.south west) -- ++(-0.55,-0.55) |- (backend.east);

\end{tikzpicture}
}
\caption{Thesis-specific compiler pipeline: the implementation moves from
source identity through static analyses, high-level transforms, recurrence
metadata, validation, and backend execution.}
\label{fig:pipeline}
\end{figure}

\section{Type Context and Pass Manager}

\pathcode{src/einlang/passes/base.py} defines the shared pass interfaces. The
\texttt{TyCtxt} object is the compiler's single context for resolver state,
source files, diagnostics, analysis results, module-system state, source
overlays, and later pass metadata. This mirrors rustc's discipline: passes
communicate through an explicit context instead of through global state.

Passes derive from \texttt{BasePass}; each exposes a \texttt{run} method that
receives the program and \texttt{TyCtxt}. The \texttt{PassManager} records pass
classes and sorts them by declared dependencies. The driver still owns a few
phase boundaries, notably parsing, name resolution, and AST-to-IR lowering.
Later pipeline steps follow the registered pass order.

\subsection{Pass Scheduling Algorithm}

The pass manager is intentionally small. Each pass declares its dependencies,
and the manager converts that dependency graph into an execution order. The
driver then runs passes over the same mutable IR and shared \texttt{TyCtxt}.

\begin{lstlisting}[style=pythoncode,caption={Dependency-ordered pass execution.}]
run_pass_pipeline(program, tcx, requested_passes):
    order = topological_sort(requested_passes, edge = pass.requires)
    for pass_class in order:
        analysis_before = tcx.analysis_table
        program = pass_class().run(program, tcx)
        assert program is not None
        tcx.mark_pass_complete(pass_class)
        validate_declared_outputs_if_enabled(program, analysis_before, tcx)
    return program

topological_sort(passes):
    temporary = set()
    permanent = set()
    out = []

    visit(pass):
        if pass in permanent:
            return
        if pass in temporary:
            report dependency cycle
        temporary.add(pass)
        for dep in pass.requires:
            visit(resolve_pass(dep))
        temporary.remove(pass)
        permanent.add(pass)
        out.append(pass)

    for pass in passes:
        visit(pass)
    return out
\end{lstlisting}

Specialization is the one place where pass scheduling becomes incremental:
type inference may discover a new specialized function, then invoke the prefix
of the pipeline needed to analyze that new body. The same \texttt{TyCtxt}
holds the shared resolver and specialization registry, so the specialized body
can be inserted back into the main program without losing identity.

The scheduling algorithm is deliberately conservative. A pass can only rely on
facts produced by its declared dependencies, and the driver records completion
after a pass has successfully returned a program. This prevents hidden ordering
contracts such as "shape analysis happens to run before type inference" from
living only in the driver's source order. The cycle check is equally important:
analysis passes often grow by consuming newly attached facts, and a dependency
cycle would otherwise turn into a confusing partial analysis rather than an
explicit compiler configuration error.

The validation hook after each pass is an implementation guardrail. Many passes
mutate IR annotations rather than creating a fresh immutable tree, so a pass can
accidentally leave stale facts behind. Capturing the analysis table before the
pass and validating declared outputs gives tests a way to assert that a pass
produced the metadata it promised. This is the practical substitute for a
heavier typed pass framework: the repository stays small, but the pipeline still
has auditable boundaries.

The algorithm is intentionally a scheduler, not a hidden optimizer. It does not
try to reorder passes based on cost or opportunistic shortcuts; it only respects
declared dependencies. That makes pass behavior easier to explain in the thesis:
range analysis is before shape analysis because shape analysis consumes ranges,
not because a particular driver happened to call it first. When a new pass is
added, the author must state what it depends on, and the topological sort either
places it legally or reports a cycle.

This matters for incremental monomorphization as well. A specialized function
created during type inference is not allowed to skip the facts that ordinary
functions receive. The driver can run the required prefix over the clone because
the pass dependencies define that prefix in reusable terms. The result is a
pipeline that can grow dynamically while preserving the same invariants as the
initial program. The small pseudo-code listing therefore captures a larger
engineering principle: compiler phases are explicit contracts, and generated IR
must satisfy the same contracts as source IR.

\textit{Concrete example.} For \texttt{let C[i,j] =
sum[k](A[i,k]*B[k,j]);}, a requested backend execution cannot run type
inference before range analysis, because type inference depends on the shape
facts derived from \texttt{i}, \texttt{j}, and \texttt{k}. The scheduler places
range analysis before shape analysis, shape analysis before type inference, and
type inference before lowering. If a later generic call creates a specialized
matrix helper during type inference, the same dependency prefix is run on the
clone before it rejoins the program. The example shows why pass order is a data
dependency, not a cosmetic list.

\section{IR Node Design}

The IR in \pathcode{src/einlang/ir/nodes.py} is a slot-based class hierarchy.
Every IR node has a source location. Expression nodes carry optional
\texttt{type\_info} and \texttt{shape\_info}. Identity-bearing nodes such as
bindings, parameters, identifiers, functions, index variables, and rest indices
carry DefIds. The visitor interfaces in \pathcode{src/einlang/ir} and
\pathcode{src/einlang/ir/scoped_visitor.py} let passes recurse through the IR
without hand-writing traversal logic for every use.

\subsection{Mutable Analysis Annotations}

Although \einlang{} source bindings are immutable, the compiler IR is not a
formal immutable SSA graph. IR nodes are slot-based Python objects, and passes
attach or update analysis fields such as \texttt{type\_info},
\texttt{shape\_info}, \texttt{range\_ir}, recurrence overrides, execution-facts
IDs, and lowered forms. This is a pragmatic compiler representation: source
values are single-assignment from the user's point of view, while compiler
metadata evolves as each pass learns more.

There is therefore no mutable-to-SSA conversion pass. The compiler instead uses
DefIds for identity and explicit IR nodes for binding, control flow, tensor
clauses, reductions, and lowered loops. If a future backend needs a CFG-style
SSA representation, that would be a new lowering target rather than a
description of the current IR.

\begin{table}[htbp]
\centering
\caption{Important IR families.}
\label{tab:ir-families}
\renewcommand{\arraystretch}{1.08}
\small
\begin{tabular}{@{}L{0.32\textwidth}L{0.56\textwidth}@{}}
\toprule
IR family & Role \\
\midrule
Expression IR & Literals, identifiers, unary and binary operations, calls, blocks, arrays, tuples, casts, member access, control flow, and builtins. \\
Tensor IR & Rectangular access, reduction expressions, Einstein bindings, Einstein clauses, and lowered tensor nodes. \\
Autodiff IR & Differential, JVP, VJP, and lazy-Jacobian request nodes. \\
Binding IR & Top-level and local bindings, functions, constants, parameters, where clauses, modules, and patterns. \\
Lowered IR & Lowered Einstein, reduction, selection, comprehension, and recurrence nodes. \\
\bottomrule
\end{tabular}
\end{table}

IR serialization in \pathcode{src/einlang/ir/serialization.py} supports
S-expression dumps. The CLI exposes this through \texttt{--dump-ir}, which is
important for debugging pass interactions.

\section{Compiler Pipeline}

The registered pipeline is ordered so that source facts are available before
passes consume them. The important passes are:

\begin{table}[p]
\centering
\caption{Main compiler passes and their responsibilities.}
\label{tab:pass-pipeline}
\renewcommand{\arraystretch}{1.08}
\small
\begin{tabular}{@{}L{0.32\textwidth}L{0.56\textwidth}@{}}
\toprule
Pass & Responsibility \\
\midrule
AST to IR lowering & Convert source AST into IR, preserving source locations, DefIds, declarations, expressions, and tensor syntax. \\
Einstein grouping & Group repeated indexed clauses for the same binding into one high-level Einstein binding. \\
Constraint classifier & Classify where-clause predicates and bindings so later range and lowering passes know which constraints are domains, predicates, or value bindings. \\
Rest preprocessing & Expand named rest patterns such as \texttt{..batch} into concrete generated index variables where possible. \\
Range analysis & Infer index ranges from explicit domains, shaped reads, reductions, comprehensions, and constraints. \\
Shape analysis & Compute shapes for expressions and bindings, including array literals, indexed bindings, reductions, offsets, and conditionals. \\
Type inference & Infer scalar precision, array ranks, function types, tuple types, differential types, and call compatibility. \\
Extremum canonicalization & Turn recognized argmax/argmin selection patterns into explicit selection IR. \\
Pre-autodiff pruning & Prune shape/rank branches before differentiation. \\
Autodiff & Expand source differential forms and collect derivative graph facts while high-level tensor IR is still visible. \\
AD request lowering & Lower JVP/VJP/lazy-Jacobian requests into ordinary IR/runtime requests. \\
AD leak check & Ensure source-level \texttt{@} artifacts do not survive past the autodiff phase. \\
Einstein lowering & Lower indexed tensor definitions and reductions into lowered loop/reduction IR. \\
Recurrence order & Infer recurrence dimensions, same-timestep dependencies, execution order hints, and storage metadata. \\
Lowered execution facts & Attach compiler-owned facts that let the backend choose vectorized, scalar, hybrid, or specialized reduction paths without re-analysis. \\
Validation passes & Validate casts, pipeline types, match exhaustiveness, and IR invariants before execution. \\
\bottomrule
\end{tabular}
\end{table}

\section{Range, Shape, and Type}

Range analysis in \pathcode{src/einlang/passes/range_analysis.py} maintains a
DefId-keyed map from index variables to range IR. It intersects multiple
constraints by taking the maximum of starts and the minimum of ends. It also
uses an implicit range detector to recover ranges from tensor reads. Shape
analysis in \pathcode{src/einlang/passes/shape_analysis.py} uses those ranges to
derive output shapes, check rectangular indexing, handle offset expressions, and
annotate expressions with shape facts. Type inference in
\pathcode{src/einlang/passes/type_inference.py} then assigns primitive,
rectangular, tuple, function, and differential types.

The separation is deliberate. Shapes are not stored inside scalar precision
types; rectangular types can mention rank and optional static dimensions, while
shape analysis carries concrete expression-level shape facts. That split keeps
type compatibility and tensor shape reasoning from collapsing into one
monolithic pass.

\subsection{Analysis Flow}

The three analyses form a dependency chain:

\begin{enumerate}
\item Range analysis discovers domains for index variables from explicit
brackets, reductions, comprehensions, shaped reads, and classified constraints.
\item Shape analysis uses those domains to assign shapes to indexed bindings,
array literals, reductions, calls, and control-flow joins.
\item Type inference uses shape and rank information to check operations,
function calls, casts, literals, tuples, and differential values.
\end{enumerate}

This order is what lets code such as matrix multiplication report an index or
shape problem at compile time rather than failing inside a NumPy call.

\begin{lstlisting}[style=einlang]
let C[i, j] = sum[k](A[i, k] * B[k, j]);
\end{lstlisting}

The compiler must infer that \texttt{i} comes from the first axis of
\texttt{A}, \texttt{j} from the second axis of \texttt{B}, and \texttt{k} from
both the second axis of \texttt{A} and the first axis of \texttt{B}. A mismatch
in those inferred \texttt{k} ranges is a source-level shape error.

This one-line program also explains why the analyses are separate. Range
analysis can discover candidate domains for \texttt{i}, \texttt{j}, and
\texttt{k} without yet deciding the element type of the multiplication. Shape
analysis can then use those ranges to infer the result rank and verify that the
contracted axis is compatible. Type inference finally checks whether the
element operations are legal. If those concerns were fused into one pass, every
diagnostic would have to disentangle domain, shape, and element-type failures
after the fact.

The line also illustrates how source structure affects error locality. If
\texttt{A} has shape \texttt{[M,K]} and \texttt{B} has shape \texttt{[L,N]},
the inconsistency is not an arbitrary runtime broadcasting error; it is the fact
that the same reduction DefId \texttt{k} is constrained by two incompatible
axes. The diagnostic can therefore point to the two reads that introduced the
conflicting domains. That is a stronger compiler story than treating the
program as a call to a library contraction and reporting whatever error the
library raises.

\subsection{Range Analysis Algorithm}

Range analysis is a DefId-indexed constraint collection pass. It stores every
candidate range for an index variable and resolves the effective range as the
intersection of those candidates. Intersection is represented in IR, not only
as integers: multiple lower bounds become \texttt{max(...)} and multiple upper
bounds become \texttt{min(...)}.

\begin{lstlisting}[style=pythoncode,caption={Range inference for indices and reductions.}]
analyze_ranges(program):
    analyzer.ranges = map DefId -> list[RangeIR]
    for function in analyzable_functions(program):
        visit_function(function)
    for statement in program.statements:
        visit(statement)
    return analyzer

visit_einstein_binding(binding):
    check all clauses have the same output rank
    for clause in binding.clauses:
        for output_index in clause.indices:
            if output_index has explicit range:
                analyzer.add_range(output_index.defid, output_index.range_ir)
                clause.variable_ranges[output_index.defid] = output_index.range_ir
        classify where-constraints as ranges, guards, or local bindings
        infer implicit ranges from shaped reads in clause.value
        copy ranges across clauses at the same output-axis position

visit_reduction(reduction):
    visit reduction.body and reduction.where constraints
    detector = ImplicitRangeDetector(visible_bindings, tcx)
    detector.infer_reduction_ranges_from_where(reduction)
    for loop_var in reduction.loop_vars:
        if loop_var has no explicit range:
            range = detector.infer_implicit_range(reduction.body, loop_var.defid)
                 or detector.infer_reduction_var_range_from_body(...)
        if range is missing:
            report source diagnostic with a suggested explicit range
        reduction.loop_var_ranges[loop_var.defid] = range

get_range(defid):
    candidates = analyzer.ranges[defid]
    start = max(candidate.start for candidate in candidates)
    end = min(candidate.end for candidate in candidates)
    return RangeIR(start, end)
\end{lstlisting}

The important detail is that the range map is keyed by DefId rather than by
printed names. This lets one function contain several independent variables
called \texttt{i}; their ranges never collide.

The algorithm is constraint collection followed by normalization. Explicit
ranges are the strongest source of information because they come directly from
the binding or reduction syntax. Implicit ranges are secondary: they are
recovered from shaped reads only when the compiler can determine which axis of
which tensor an index variable inhabits. Where-clause classification happens in
the same pass boundary because predicates and value bindings can look similar
at the surface but lower to very different IR. A predicate becomes a guard; a
range domain becomes a loop bound; a value binding becomes a local computation.

Representing intersections as IR rather than as immediate integers is what
lets symbolic programs remain analyzable. If one use gives \texttt{0..N} and
another gives \texttt{0..M}, the effective range can be represented as
\texttt{0..min(N,M)} even when neither extent is known yet. Later shape
analysis may simplify that expression after more dimensions are known. The
failure path is also part of the design: when a reduction variable cannot be
inferred, the compiler reports a source diagnostic asking for an explicit
range instead of silently choosing a full tensor extent.

The algorithm also explains why implicit inference is deliberately limited. A
single index occurrence in a shaped read is informative only when the compiler
can pair that occurrence with an array axis. Arithmetic expressions such as
\texttt{i + r} contribute bounds only after the participating indices already
have candidate ranges. This avoids circular reasoning: the compiler may refine a
known range with an offset constraint, but it should not invent a domain from an
unanchored arithmetic expression. The user-facing consequence is predictable:
write explicit ranges when shape reads do not determine the domain.

Range analysis therefore serves two audiences. Shape analysis consumes the
normalized ranges to allocate tensor extents, while diagnostics consume the
candidate list to explain where each bound came from. Keeping the candidates,
rather than only the final \texttt{max}/\texttt{min} expression, lets the
compiler distinguish an explicit user range from an inferred axis range. That is
useful when a future warning or optimization wants to explain why a loop is
smaller than a source dimension or why two symbolic bounds could not be proven
compatible.

\textit{Concrete example.} If \texttt{A} has shape \texttt{[M,K]} and
\texttt{B} has shape \texttt{[K,N]}, the reduction variable \texttt{k} in
\texttt{A[i,k] * B[k,j]} receives two candidate ranges, both \texttt{0..K}; the
effective range remains \texttt{0..K}. If the second tensor instead has shape
\texttt{[L,N]}, the range map records candidates \texttt{0..K} and
\texttt{0..L} for the same DefId and normalizes them to
\texttt{0..min(K,L)} while preserving the original sources of the constraints
for diagnostics. If \texttt{k} appears only in \texttt{k + 1} with no shaped
read or explicit \texttt{k in 0..K}, the pass reports that no domain can be
inferred.

\subsection{Shape Analysis Algorithm}

Shape analysis consumes ranges and writes two maps: expression-to-shape and
DefId-to-shape. The DefId map lets an identifier recover the shape of the value
it names even when the identifier node itself has not been seen before.

\begin{lstlisting}[style=pythoncode,caption={Einstein shape inference.}]
infer_einstein_shape(binding):
    clauses = binding.clauses
    rank = first non-empty clause index count
    clause_shapes = []

    for clause in clauses:
        axis_extents = []
        for index in clause.indices:
            if index is a literal:
                axis_extents.append(index.value + 1)
            else if clause.variable_ranges[index.defid] has static end:
                axis_extents.append(end_value)
            else:
                axis_extents.append(infer_extent_from_array_reads(index, clause.value))
        clause_shapes.append(tuple(axis_extents))

    outer_shape = dimensionwise_max(clause_shapes)

    value_shapes = []
    for clause in clauses:
        value_shape = infer_value_shape(clause.value)
        if value_shape is known:
            value_shapes.append(value_shape)

    if value_shapes is empty:
        return outer_shape
    if all value_shapes are identical:
        return outer_shape + value_shapes[0]
    report or defer shape mismatch
\end{lstlisting}

Dependent ranges are resolved opportunistically. Range analysis may leave an
axis as \texttt{0..x.shape[dim]}. Once shape analysis knows the shape of
\texttt{x}, it rewrites the end of that range to a literal dimension. That
rewrite lets later lowering and backend planning see a concrete loop bound.

The clause loop in the algorithm has to account for piecewise tensor
definitions. A literal index such as \texttt{0} contributes a covered point, not
an output loop variable, while an index variable contributes an axis whose
extent must be inferred or read from the explicit domain. Taking the
dimensionwise maximum across clause shapes gives the allocation shape for a
rectangular backend while preserving the fact that not every clause writes the
whole tensor. This is why boundary clauses and interior clauses can coexist in
one binding.

The value-shape portion handles tensors whose elements are themselves shaped
values. An indexed binding may produce scalars, vectors, or higher-rank tensor
elements. The output-index shape is therefore only the outer part of the result;
the body shape is appended when all clauses agree. If body shapes disagree, the
compiler can report a mismatch at the clause level rather than allowing a
ragged rectangular tensor to reach the backend. This check is a direct example
of source structure improving diagnostics.

The algorithm separates allocation shape from value shape because indexed
definitions can describe tensor-valued elements. For scalar-valued clauses the
outer shape is the whole result. For vector-valued or matrix-valued clause
bodies, the output axes describe where each element lives and the body shape
describes what each element contains. This distinction prevents a common
implementation shortcut: flattening everything into one rank and later guessing
which axes belonged to the source. \einlang{} keeps the source axes and element
axes explicit until lowering.

The dimensionwise maximum is conservative but useful for piecewise definitions.
A base clause that writes \texttt{0} and an interior clause that writes
\texttt{1..N} must agree on one rectangular result allocation, even though their
coverage sets differ. Shape analysis computes that allocation while leaving
coverage and guard information on the clauses. The backend can then allocate the
right container without pretending that every clause covers every point. This is
one of the places where the high-level IR keeps more information than a simple
loop nest would.

\textit{Concrete example.} The three-clause Fibonacci binding writes literal
points \texttt{fib[0]} and \texttt{fib[1]} plus a ranged clause
\texttt{fib[n in 2..N]}. Shape analysis treats the literal clauses as covered
points and the ranged clause as the axis extent, so the allocation shape is the
single sequence axis rather than three unrelated scalar outputs. For a binding
such as \texttt{let rows[i] = [x[i], y[i]]}, the outer shape comes from
\texttt{i}, while the body contributes an element shape of length two; the
result is a rectangular two-dimensional value whose axes remain distinguishable
for lowering and diagnostics.

\subsection{Type Inference Algorithm}

The current type pass is a checking and propagation pass, not full Hindley-
Milner inference. It walks IR in scoped order, assigns types to nodes, checks
operator compatibility, registers parameter and pattern bindings, and asks the
monomorphization service for specialized copies of generic functions when call
argument types become informative.

\begin{lstlisting}[style=pythoncode,caption={Type checking with specialization feedback.}]
infer_types(program):
    inferencer = TypeInferencer(tcx)
    program.accept(inferencer)

    while monomorphization_service.has_pending_functions():
        pending = monomorphization_service.take_pending_functions()
        for specialized in pending:
            insert specialized into program and function_ir_map
            if specialized.return_type is unknown:
                specialized.accept(inferencer)

    rewrite generic call sites to specialized DefIds where available
    return program

visit_call(call):
    infer argument types
    callee = resolve function_defid
    if callee is generic and argument types expose precision or rank:
        spec = monomorphization_service.incremental_monomorphize(
            call, argument_types, current_pass="type",
            required_passes=["einstein_grouping", "constraint_classifier",
                             "rest_pattern", "range", "shape", "type"])
        if spec exists:
            call.function_defid = spec.defid
            return spec.return_type
    check argument compatibility against callee parameters
    return callee.return_type
\end{lstlisting}

The type pass is intentionally described as feedback-driven rather than purely
bottom-up. A call can reveal enough information to specialize a generic
function, and that specialization may need earlier passes before its return
type is available. The pass therefore alternates between ordinary inference and
processing pending specialized functions. This keeps the public language
generic while still giving later passes concrete DefIds and more precise
parameter types.

The algorithm also draws the current boundary of inference. It checks and
propagates types that the implementation knows how to represent; it does not
attempt principal-type inference for the whole language. That boundary is a
good fit for the compiler because most hard tensor questions depend on shape
and range metadata rather than on polymorphic type schemes. When argument
precision or rank is informative, the monomorphization service clones a body
and lets the established pass pipeline analyze it. When the call is too
ambiguous, the compiler leaves it generic until more information is available
or reports a compatibility error.

The feedback loop also controls when generic functions become ordinary analyzed
functions. A call with only dynamic-rank information may not justify a clone; a
call with known precision or rank can produce a body that shape analysis and
lowering can reason about more concretely. The pass records that choice through
DefIds rather than through a runtime tag, so after rewriting the call site the
specialized function is just another function in the program. This keeps later
passes simple: they do not need to know whether a body was written by the user
or generated by monomorphization.

The current algorithm favors explicit checks over global inference. That is a
reasonable boundary for \einlang{} because most examples are numerical kernels
whose hard questions are about axes, shapes, and derivative structure. Full
principal-type inference would add complexity without replacing the need for
range and shape analysis. The implemented pass therefore spends its precision
budget where the language needs it most: typed calls, scalar operations,
tensor ranks, and the creation of specialized bodies that downstream passes can
validate.

\textit{Concrete example.} A generic helper \texttt{fn twice(x) \{ x + x \}}
can be called with a scalar \texttt{f32}, a vector \texttt{[f32]}, and a matrix
\texttt{[f32; *, *]}. Type inference checks the addition in each context and,
when rank or precision is informative, asks monomorphization for a specialized
body. The scalar call returns an \texttt{f32}; the vector call returns a vector
with the argument's element precision; and the matrix call keeps rank-two shape
facts available for later indexed uses. A call with an unsupported pair such as
a tuple and a tensor fails during type checking before lowering.

\section{Monomorphization}

Generic functions are specialized by
\pathcode{src/einlang/analysis/monomorphization_service.py}. The service keys
specializations by generic DefId and normalized parameter types. It can create
full specializations when both precision and rank are known, and partial
specializations when only precision or only rank is available. Newly specialized
functions are analyzed incrementally by the relevant passes and then inserted
into the program before tree shaking.

This matters for the standard library: many library functions are written once
but called at different ranks or precisions. Monomorphization lets the compiler
keep the source surface generic while giving later analysis and backend code
more concrete IR.

The service is incremental because specialization can be discovered while type
inference is already walking a program. A call site can request a specialized
copy, the service can clone and register that copy with a new DefId, and the
relevant analysis passes can run over the new body. After specialization, call
sites are rewritten to the specialized DefId so the backend does not perform
generic dispatch at runtime.

\subsection{Specialization Algorithm}

The specialization key is the generic function DefId plus normalized parameter
types. Normalization preserves the facts needed by the compiler while avoiding
one copy per concrete tensor extent. For rectangular tensors, concrete extents
are replaced by \texttt{None} at each known rank position.

\begin{lstlisting}[style=pythoncode,caption={Incremental monomorphization.}]
incremental_monomorphize(call, arg_types, required_passes):
    if call target is not generic:
        return None
    if arg_types expose neither precision nor rank:
        return None
    if specialization count for target exceeds safety limit:
        return None

    key = Instance(call.function_defid, normalize(arg_types))
    if key in registry:
        spec = registry[key]
        run_missing_required_passes(spec, required_passes)
        call.function_defid = spec.defid
        return spec

    if target is already being monomorphized:
        return None

    clone = deep_copy(generic_function)
    clone.defid = resolver.allocate_for_item(module_path, specialized_name)
    for parameter, arg_type in zip(clone.parameters, arg_types):
        parameter.param_type = dynamic_shape_type(arg_type)
    copy custom differential rule from generic function if needed

    registry[key] = clone
    function_ir_map[clone.defid] = clone
    pending_specialized_functions.append(clone)
    run_required_prefix_passes(clone, required_passes)
    call.function_defid = clone.defid
    return clone
\end{lstlisting}

The service also performs a small dead-code-elimination pass over the cloned
body before analysis. Conditions such as \texttt{len(x.shape) == 2} can become
compile-time constants after parameter ranks are known, so unreachable branches
are removed before range and shape analysis see them.

The specialization key deliberately ignores concrete extents for rectangular
tensors while preserving rank and element facts. This prevents a function from
being cloned once for every input length, but still lets the compiler
distinguish scalar, vector, matrix, and dynamic-rank behavior. The safety limit
and re-entrancy check are defensive: generic code can be recursive or can
produce many call shapes, so specialization must not turn a small program into
an unbounded clone graph.

Copying custom differential rules along with the function body is also part of
semantic preservation. A generic function and its specialized copy should agree
about how autodiff treats the call. Once the specialized DefId is installed,
the call site no longer needs runtime generic dispatch; range, shape, type,
autodiff, and backend passes can reason about a concrete function body. The
analysis value is therefore not only performance. Specialization makes later
compiler facts more precise while keeping source code reusable.

The clone step must also preserve identity hygiene. Parameters in the cloned
function receive specialized type information, but the function itself receives
a new DefId so uses of the generic body and uses of the specialized body do not
collide. Any local DefIds inside the clone are treated as belonging to that
clone's analysis context. This is why specialization is routed through the same
resolver and function map as source functions rather than being stored as an
opaque backend cache.

The partial-specialization policy is a practical compromise. Preserving rank but
not concrete extent lets a matrix helper specialize once for all two-dimensional
inputs of the same element family. Preserving precision lets numeric builtins
avoid repeatedly checking scalar families at runtime. At the same time, ignoring
extent prevents common workloads from producing one clone per batch size. The
algorithm is therefore a semantic precision tool with an explicit blow-up
guard, not an unconditional inliner.

\textit{Concrete example.} Suppose a stdlib function branches on rank:
\texttt{if len(x.shape) == 2 \{ mat\_case(x) \} else \{ vec\_case(x) \}}. When
the call site passes a matrix, specialization normalizes the argument type to
"rank two, element family known, extents dynamic" and clones a rank-two body.
The dead branch for non-matrices can be pruned before range and shape analysis,
so the backend never sees a runtime rank test for that call. Passing matrices
of shapes \texttt{32x32} and \texttt{64x64} reuses the same specialization
because concrete extents are not part of the key.

\section{Diagnostics and Validation}

Diagnostics are implemented in \pathcode{src/einlang/shared/errors.py}. Errors
carry source locations, optional codes, labels, notes, help messages, and
multi-span labels. The formatter renders rustc-style snippets when the source
text is available. Compiler-driver exception handling turns uncaught pass
exceptions into diagnostics tied to either the current IR span or the original
source file.

Validation is split across focused passes. Cast validation checks explicit
numeric conversions. Pipeline validation checks pipeline expression types.
Exhaustiveness validates match coverage. IR validation checks post-lowering
invariants, including the requirement that semantic identity be DefId-based.
Tree shaking removes unused functions and modules after all analyses have had a
chance to run.

\subsection{Validation Algorithms}

Validation passes are deliberately narrow. For example, exhaustiveness checking
does not try to prove every numeric interval; it implements the cases currently
supported by the pattern language and reports precise missing cases when it
can.

\begin{lstlisting}[style=pythoncode,caption={Current match exhaustiveness check.}]
check_match(match_expr):
    type_name = primitive_name(match_expr.scrutinee.type_info)
    patterns = flatten_or_patterns(arm.pattern for arm in match_expr.arms)

    if any(pattern is wildcard or identifier binding for pattern in patterns):
        return exhaustive

    if type_name is bool:
        has_true = any(pattern is literal True for pattern in patterns)
        has_false = any(pattern is literal False for pattern in patterns)
        if not has_true or not has_false:
            report E0004 with missing true/false labels
        return has_true and has_false

    if type_name is integer:
        report E0004 unless a wildcard-like pattern exists

    otherwise:
        conservatively report non-exhaustive
\end{lstlisting}

The same design applies to post-lowering validation: the compiler checks the
invariants that later code assumes, such as non-null DefIds on semantic
identifiers, expanded rest indices before range analysis, no source-level
autodiff nodes after request lowering, and compiler-owned recurrence and
vectorization metadata before the NumPy backend executes lowered tensor IR.

The exhaustiveness algorithm is intentionally modest, but that modesty is a
feature of the current implementation rather than a gap in the thesis story.
For booleans, the compiler can give exact missing-case diagnostics because the
domain has two values. For integers, the current pass requires a wildcard-like
case rather than proving interval coverage. For other types, it reports
conservatively unless a catch-all pattern is present. This avoids pretending to
prove properties that the pattern language and type information do not yet
support.

The broader validation lesson is that each pass enforces the invariants needed
by the next implementation boundary. Exhaustiveness protects match execution,
cast validation protects numeric conversion, the autodiff leak check protects
the backend from source-level \texttt{@} nodes, and IR validation protects
DefId-based lookup. The code snippet is therefore one example of a repeated
pattern: narrow validators produce clearer errors and make backend assumptions
explicit.

The analysis value of the validator is that it documents what the compiler does
not yet prove. Boolean exhaustiveness is exact, integer exhaustiveness is
conservative, and richer pattern domains are reported conservatively unless a
catch-all exists. That conservatism is preferable to a validator that accepts a
match by accident and lets the interpreter discover a missing arm later. The
validator's result is therefore a phase certificate: execution may assume that
the currently supported pattern fragment has been checked.

The same philosophy applies to leak checks and IR validation. They are not
optimizations; they are boundary tests between compiler phases. If source-level
autodiff reaches the backend, the backend would have to decide what \texttt{@}
means after the phase designed for that meaning has already passed. If an
identifier lacks a DefId after resolution, runtime lookup would have to fall
back to strings. The validators reject these states so backend code can stay
small and direct.

\textit{Concrete example.} The expression \texttt{match flag \{ true => 1 \}}
is rejected for a boolean scrutinee because the pass can name the missing
\texttt{false} case exactly. The expression \texttt{match n \{ 0 => 1 \}} is
also rejected for an integer scrutinee, but for a different reason: the current
implementation does not prove full integer interval coverage, so it requires a
wildcard-like arm. After autodiff lowering, a remaining \texttt{@loss/@W} node
is likewise rejected as a phase leak rather than being passed to the backend.
